
\subsubsection{2.1 Pre-Training Corpus Curation and Model Generalization}

The empirical literature on corpus curation and model performance is substantial but consistently confounded by architecture differences or training duration mismatches. Biderman et al. [2023] demonstrate via the Pythia suite that deduplication yields small but consistent benchmark improvements, including a modest HellaSwag gain of approximately 0.01 absolute. This within-architecture, within-corpus comparison isolates deduplication as a single intervention; it does not address multi-stage quality curation, and the small HellaSwag gain is consistent with near-saturation behavior at 300 billion tokens. Penedo et al. [2023] show that web-filtered RefinedWeb data outperforms unfiltered training data for Falcon models; however, the comparison involves different model architectures and focuses on average performance rather than generalization balance ratios, leaving the architecture confound unresolved.

Groeneveld et al. [2024] document OLMo-7B's design and report competitive MMLU performance at full training (approximately 2 trillion tokens), attributing part of this to Dolma's academic content inclusion. This comparison, however, is not matched by training scale; comparing OLMo at 2 trillion tokens to Pythia at 300 billion tokens conflates corpus quality with training duration. The present work tests the same corpus pairing at a matched intermediate scale.

Muennighoff et al. [2023] establish a data quality $\times$ training scale interaction, finding that higher-quality data benefits may compound with additional training tokens. A null result at 300 billion tokens is consistent with this interaction: the Dolma advantage may emerge at later training stages without implying that curation is ineffective overall.

Soldaini et al. [2024] document the Dolma corpus in detail, including its multi-stage curation pipeline, domain composition, and deduplication methodology. This documentation provided the primary basis for the hypothesis that Dolma would produce superior generalization balance relative to The Pile.

\subsubsection{2.2 Generalization Evaluation Metrics}

MMLU \citep{Hendrycks2021MMLU} is a 57-subject multiple-choice benchmark spanning STEM, humanities, social sciences, and professional domains, designed to measure multi-task language understanding. HellaSwag \citep{Zellers2019HellaSwag} is a commonsense completion benchmark derived from ActivityNet and WikiHow, designed to measure web-sourced procedural and situational knowledge. ARC-Easy and ARC-Challenge \citep{Clark2018ARC} are science question benchmarks at elementary and challenge levels.

The use of performance ratios as generalization balance metrics has precedent in out-of-distribution generalization research, where in-distribution/out-of-distribution performance gaps characterize model behavior [Miller et al., 2021; Koh et al., 2021]. However, ratio metrics require that both numerator and denominator remain discriminative at the comparison point. The present study demonstrates that this requirement is violated when the denominator task saturates at the target training scale.

\subsubsection{2.3 Data Attribution and Quality Proxy Methods}

Influence functions \citep{Koh2017Influence}, TRAK \citep{Park2023TRAK}, and TracIn \citep{Pruthi2020TracIn} provide tools for attributing model behavior to specific training examples. These methods could in principle quantify the contribution of high-quality versus low-quality data subsets to generalization performance but have not been applied to the corpus-quality $\times$ generalization-balance question at 300-billion-token pre-training scale. Such attribution analyses remain a promising direction but require architecture-matched designs to avoid confounding attribution results with architecture effects.

\subsubsection{2.4 Positioning}

This work differs from prior literature in three respects. First, a matched training scale (approximately 300 billion tokens, both models within 0.5\% of target) eliminates training duration as a confound. Second, the focus is on generalization balance --- the ratio between knowledge-intensive and commonsense task performance --- rather than absolute benchmark scores. Third, a scope-qualified null result is reported with explicit characterization of the conditions under which the comparison breaks down, rather than drawing causal conclusions from an architecturally confounded design.

